\section{What current research adds}\label{ch14:sec:research}
The question of this chapter is a well-known research problem. The \emph{circulant Hadamard conjecture}, usually attributed to H.~J.~Ryser, states that there is no circulant Hadamard matrix of order greater than $4$. In the language of this chapter: apart from the lengths $1$ and $4$, no sign sequence has all of its periodic correlations at nonzero shifts equal to zero. The conjecture has been open since the 1960s. Proposition~\ref{thm:circulantnecessary} shows that a counterexample would need an order $4u^2$ with $u\ge2$. Over the decades, researchers have proved much stronger restrictions on the possible orders, but none of these results rules out all of them.

Barker sequences have a related history. Since the early 1960s it has been known that a Barker sequence of odd length has length at most $13$. In our table, the odd lengths that have Barker sequences are exactly $1$, $3$, $5$, $7$, $11$, and $13$. For even lengths, the Barker question turns into a question about circulant Hadamard matrices. It has long been known that when $n>2$ is even, every Barker sequence of length $n$ also has $C_s=0$ for every shift $s$ with $1\le s\le n-1$, so its circulant matrix is a Hadamard matrix. We do not prove this here, but you can see one instance of it in the table. The eight Barker sequences of length four are $(1,1,1,-1)$, its three other rotations, and the negatives of these four. They are also the eight sequences counted in the last column. As a consequence, a Barker sequence of even length greater than $4$ would give a circulant Hadamard matrix of order greater than $4$, and whether such a sequence exists was likewise an open question.

In September 2026, a preprint claiming a proof of the circulant Hadamard conjecture was released in OpenAI's public mathematics repository, together with a formalization in Lean~\cite{oai179}. Lean is a \emph{proof assistant}: a computer program that checks every step of a proof written in its own precise language. A \emph{formalization} of a theorem is a version of its statement and proof written in such a language, so that the program can check it. Chapters~\ref{ch:lean1} and~\ref{ch:lean2} introduce Lean.

According to the repository's description, the formalized statement says that a real circulant Hadamard matrix of positive order $n$ exists exactly when $n=1$ or $n=4$. The word ``real'' separates the matrices of this chapter, whose entries are the real numbers $1$ and $-1$, from a related family with complex entries that researchers also study. A second formalized statement, described there as a consequence, says that a sign sequence of positive even length whose aperiodic correlations at all nonzero shifts have absolute value at most one must have length $2$ or $4$. In our terms, every Barker sequence of even length has length $2$ or $4$. The description also notes that the preprint's treatment of odd lengths is not part of the formalized statements.

These claims come from an outside source, and nothing in this chapter establishes them. Our experiment stops at length $13$, and Proposition~\ref{thm:circulantnecessary} still permits the orders $16,36,64,\ldots$. What the chapter does give you is the ability to read the claims precisely. You know what a circulant matrix is, how the orthogonality of its rows translates into periodic correlations, why even-length Barker sequences belong to the same story, and exactly which orders an elementary argument leaves open. Ruling out the orders $4u^2$ with $u\ge2$ requires far more mathematics than this course covers. When you meet a summary of a result like this, whether in a news report or from an AI assistant, first check which statement was proved: which matrices, which kind of correlation, which lengths. Chapter~\ref{ch:bridge} returns to this example and describes how to begin reading a proof whose background you do not yet have.
